\documentclass[a4paper,12pt]{book} % nie: report!


\usepackage[T1,plmath]{polski} % lepiej to zamiast babel!
\usepackage[utf8]{inputenc} % w razie kłopotów spróbować: \usepackage[utf8x]{inputenc}
\usepackage{fancyhdr} % nagłówki i stopki
\usepackage{indentfirst} % WAŻNE, MA BYĆ!
\usepackage[pdftex]{graphicx} % to do wstawiania rysunków
\usepackage{amsfonts}
\usepackage{amsmath}
\usepackage{amssymb}
\usepackage{amsthm}
\usepackage[pdftex,
            left=1.1in,right=1.1in,
            top=1.1in,bottom=1.1in]{geometry} % marginesy
\usepackage{float}
\usepackage[font=small,labelfont=bf]{caption}
\usepackage{xcolor} % do widocznych znaczników cytowań
\usepackage{cite}

\usepackage[colorlinks=true]{hyperref} % odnośniki interaktywne w PDFie
\hypersetup{allcolors=blue}
\usepackage{listings}
\lstset{
    basicstyle=\footnotesize\tt,
    numbers=left,
    numberstyle=\tiny,
    frame=tb,
    tabsize=4,
    columns=fixed,
    showstringspaces=false,
    showtabs=false,
    keepspaces,
    captionpos=t,
    commentstyle=\color{red},
    keywordstyle=\color{blue}
}
\newfloat{lstfloat}{htbp}{lolst}[chapter]
\floatname{lstfloat}{Listing}
\def\lstfloatautorefname{Listing}

% Widoczny, łatwy do wyszukania znacznik miejsca na przyszły odnośnik
% bibliograficzny. Nie tworzy żadnego wpisu w bibliografii.
\newcommand{\todocite}[1]{\textsuperscript{\textcolor{red}{[CYT:~#1]}}}

% skrót do nazw elementów kodu
\newcommand{\kod}[1]{\texttt{#1}}


% definicje nagłówków i stopek
\pagestyle{fancy}
\renewcommand{\chaptermark}[1]{\markboth{#1}{}}
\renewcommand{\sectionmark}[1]{\markright{\thesection\ #1}}
\fancyhf{}
\fancyhead[LE,RO]{\footnotesize\bfseries\thepage}
\fancyhead[LO]{\footnotesize\rightmark}
\fancyhead[RE]{\footnotesize\leftmark}
\renewcommand{\headrulewidth}{0.5pt}
\renewcommand{\footrulewidth}{0pt}
\addtolength{\headheight}{1.5pt}
\fancypagestyle{plain}{\fancyhead{}\cfoot{\footnotesize\bfseries\thepage}\renewcommand{\headrulewidth}{0pt}}


% interlinia
\linespread{1.25}


\begin{document}

\tableofcontents{}

\chapter{Opis implementacji}
\label{ch:opis-implementacji}

W~niniejszym rozdziale opisano wszystkie autorskie implementacje czterech jąder
obliczeniowych algebry liniowej zaimplementowanych w~ramach pracy: iloczynu
skalarnego (\kod{dot}), mnożenia macierzy przez wektor (\kod{gemv}), mnożenia
macierzy przez macierz (\kod{gemm}) oraz splotu dwuwymiarowego (\kod{conv}).
Każde jądro występuje w~wielu wariantach, uporządkowanych od implementacji
naiwnej (referencyjnej) do w~pełni zoptymalizowanej; wariantom autorskim
towarzyszą wywołania bibliotek wzorcowych (OpenBLAS, BLIS, libxsmm) służące za
punkt odniesienia. Dla każdej implementacji podano zwięzły opis zastosowanej
techniki oraz --- co jest celem tego rozdziału --- wypunktowane uzasadnienie
każdej podjętej decyzji projektowej (dobór rozmiarów kafelków, faktorów
blokowania, alokacji rejestrów, układu pakowania danych, strategii
zrównoleglenia, warunków przełączania ścieżek czy decyzji o~użyciu wstawek
asemblerowych).

Wszystkie konkretne stałe (\kod{MR\_Z}, \kod{KC\_Z}, \kod{CB} itd.) cytowane są
bezpośrednio z~kodu źródłowego. Znaczniki postaci [CYT:~klucz] oznaczają miejsca,
w~których w~ostatecznej wersji pracy powinny pojawić się odnośniki
bibliograficzne; same wpisy bibliograficzne nie są tu jeszcze zamieszczane.

\section{Środowisko sprzętowe i~wspólne założenia}
\label{sec:srodowisko}

Wszystkie jądra projektowano pod kątem jednego, konkretnego mikroprocesora ---
AMD~Ryzen~5~5600 (mikroarchitektura Zen~3). Decyzje projektowe niższego poziomu
wynikają wprost z~właściwości tej architektury:

\begin{itemize}
\item \textbf{Rejestry wektorowe.} Zen~3 udostępnia 16~architektonicznych
  rejestrów \kod{YMM} o~szerokości 256~bitów (4~liczby \kod{double} lub
  8~liczb \kod{float})\todocite{agner-microarch}. Liczba ta jest twardym
  ograniczeniem decydującym o~maksymalnym rozmiarze kafelka mikrojądra ---
  wykorzystywaną wielokrotnie w~doborze kształtów $4\times12$ (\kod{gemm}) oraz
  $6\times16$ (\kod{conv}).
\item \textbf{Jednostki FMA.} Architektura wykonuje instrukcję
  \kod{vfmadd}\,\ldots\ z~przepustowością dwóch operacji na takt i~opóźnieniem
  rzędu czterech taktów\todocite{agner-instr}. Aby ukryć to opóźnienie,
  mikrojądra utrzymują kilka niezależnych łańcuchów akumulacji (np. 8~akumulatorów
  w~iloczynie skalarnym, 12~w~\kod{gemm}).
\item \textbf{Hierarchia pamięci podręcznej.} L1d~=~32~KiB, L2~=~512~KiB,
  L3~=~32~MiB (współdzielona w~obrębie CCX)\todocite{amd-sog}. Wartości te
  wyznaczają faktory blokowania \kod{KC}, \kod{MC} i~\kod{NC}, tak aby kolejne
  panele danych rezydowały na właściwych poziomach hierarchii.
\item \textbf{Przepustowość pamięci.} Zmierzone przepustowości efektywne
  (DDR4-3200 oraz poszczególne poziomy cache) stanowiące górne ograniczenie
  modelu \emph{roofline} przejęto z~wcześniejszych pomiarów autora dotyczących
  jądra \kod{gemv}\todocite{gemv-paper}.
\end{itemize}

Wspólnym elementem niemal wszystkich jąder jest \emph{redukcja pozioma}
rejestru wektorowego do skalara (funkcja \kod{hsum\_pd}) oraz obsługa skalarnych
,,ogonów'' (elementów niemieszczących się w~pełnym kafelku wektorowym). Atrybuty
\kod{\_\_attribute\_\_((hot))} oraz \kod{always\_inline} stosowane są do
wymuszenia agresywnego inliningu mikrojąder.

\section{Iloczyn skalarny (\kod{dot})}
\label{sec:dot}

Iloczyn skalarny dwóch wektorów $a,b\in\mathbb{R}^{n}$,
\begin{equation}
s = \sum_{i=0}^{n-1} a_i\,b_i,
\label{eq:dot}
\end{equation}
jest operacją o~bardzo niskiej intensywności arytmetycznej (jedna operacja FMA
na dwa odczyty z~pamięci), więc o~jego wydajności decyduje przepustowość pamięci
oraz zdolność do ukrycia opóźnienia FMA. Kolejne warianty: \kod{naive}
$\rightarrow$ \kod{simd} $\rightarrow$ \kod{simd\_multiacc} $\rightarrow$ \kod{omp}.

\subsection{Wariant naiwny (\kod{dot\_naive})}
Pojedyncza pętla skalarna sumująca $a_i b_i$ do jednego akumulatora. Stanowi
punkt odniesienia mierzący narzut samego kompilatora i~opóźnienie pamięci bez
jawnej wektoryzacji.

\subsection{Wariant SIMD (\kod{dot\_simd})}
Wektoryzacja AVX2: pętla przetwarza po 4~liczby \kod{double} na iterację
(\kod{\_mm256\_fmadd\_pd}) do jednego akumulatora \kod{\_\_m256d}, a~na końcu
wykonuje jedną redukcję poziomą.
\begin{itemize}
\item \textbf{Pojedynczy akumulator} minimalizuje presję na rejestry, ale tworzy
  jeden łańcuch zależności FMA --- przy krótkich wektorach opóźnienie FMA nie
  jest ukryte, co jest świadomie zaakceptowanym ograniczeniem tego wariantu.
\item \textbf{Skalarny ogon} obsługuje elementy spoza wielokrotności~4.
\end{itemize}

\subsection{Wariant z~wieloma akumulatorami (\kod{dot\_simd\_multiacc})}
Kluczowy wariant zoptymalizowany: 8~niezależnych akumulatorów
(\kod{s0}\,\ldots\,\kod{s7}), pętla rozwinięta o~czynnik 32 (8~akumulatorów
$\times$~4~liczby), a~następnie redukcja drzewiasta. Implementację przedstawia
Listing~\ref{lst:dot-multiacc}.
\begin{itemize}
\item \textbf{Osiem łańcuchów akumulacji} całkowicie ukrywa czterotaktowe
  opóźnienie FMA: przy przepustowości dwóch FMA na takt potrzeba co najmniej
  ośmiu niezależnych operacji ,,w~locie''\todocite{agner-instr}.
\item \textbf{Rozwinięcie o~32} dobrano tak, aby 8~akumulatorów
  (\kod{ymm}-y wyniku) wraz z~rejestrami odczytu wypełniało plik rejestrów bez
  spillowania.
\item \textbf{Redukcja drzewiasta} ($t_0{=}s_0{+}s_1,\ldots$, następnie sumy
  częściowe) skraca ścieżkę krytyczną redukcji końcowej względem łańcucha
  liniowego.
\end{itemize}

\begin{lstfloat}[htbp]
\lstset{language=C}
\begin{lstlisting}[caption={Jądro iloczynu skalarnego z~ośmioma niezależnymi akumulatorami (\kod{dot\_simd\_multiacc.c}).}, label={lst:dot-multiacc}]
__m256d s0=_mm256_setzero_pd(), s1=_mm256_setzero_pd(), /* ... */ s7=_mm256_setzero_pd();
for (; i + 31 < n; i += 32) {
    s0 = _mm256_fmadd_pd(_mm256_loadu_pd(&a[i+ 0]), _mm256_loadu_pd(&b[i+ 0]), s0);
    s1 = _mm256_fmadd_pd(_mm256_loadu_pd(&a[i+ 4]), _mm256_loadu_pd(&b[i+ 4]), s1);
    /* ... s2..s6 ... */
    s7 = _mm256_fmadd_pd(_mm256_loadu_pd(&a[i+28]), _mm256_loadu_pd(&b[i+28]), s7);
}
__m256d t0=_mm256_add_pd(s0,s1), t1=_mm256_add_pd(s2,s3);
__m256d t2=_mm256_add_pd(s4,s5), t3=_mm256_add_pd(s6,s7);
__m256d acc=_mm256_add_pd(_mm256_add_pd(t0,t1), _mm256_add_pd(t2,t3));
\end{lstlisting}
\end{lstfloat}

\subsection{Wariant wielowątkowy (\kod{dot\_omp})}
Zrównoleglenie przez \kod{\#pragma omp parallel reduction(+:total)}: wektor jest
dzielony na równe fragmenty, a~każdy wątek wywołuje na swoim fragmencie
zoptymalizowane jądro \kod{dot\_simd\_multiacc}.
\begin{itemize}
\item \textbf{Ponowne użycie jądra skalarnego SIMD} zamiast osobnego kodu
  wielowątkowego --- każdy wątek korzysta z~w~pełni zoptymalizowanej ścieżki.
\item \textbf{Redukcja OpenMP} sumuje wyniki częściowe wątków bez jawnej
  synchronizacji.
\end{itemize}

\section{Mnożenie macierz--wektor (\kod{gemv})}
\label{sec:gemv}

Operacja \kod{gemv} realizuje
\begin{equation}
y \leftarrow \alpha\,A\,x + \beta\,y,\qquad A\in\mathbb{R}^{m\times n},
\label{eq:gemv}
\end{equation}
w~układzie wierszowym (\emph{row-major}). Podobnie jak iloczyn skalarny jest
ograniczona przepustowością pamięci, dlatego optymalizacje koncentrują się na
ponownym wykorzystaniu wczytanego wektora $x$ przez wiele wierszy jednocześnie
oraz na dopasowaniu rozmiarów bloków do hierarchii cache.

\subsection{Wariant naiwny (\kod{gemv\_naive})}
Podwójna pętla: dla każdego wiersza sumowane są iloczyny $A_{ij}x_j$, po czym
zapisywane jest $y_i = \alpha\,\mathrm{sum} + \beta\,y_i$. Implementacja
referencyjna.

\subsection{Wariant SIMD (\kod{gemv\_simd})}
Na wiersz przypadają dwa akumulatory \kod{YMM} przetwarzające 8~kolumn na
iterację (dwa bloki po 4). Obsługa $\beta$: \kod{memset} dla $\beta=0$ lub
mnożenie skalarne dla $\beta\neq1$.
\begin{itemize}
\item \textbf{Dwa akumulatory} częściowo rozdzielają zależności w~ramach jednego
  wiersza.
\item \textbf{Redukcja pozioma raz na wiersz}, po zakończeniu pętli kolumn.
\end{itemize}

\subsection{Wariant z~prefetchingiem (\kod{gemv\_simd\_prefetch})}
Rozszerza \kod{gemv\_simd} o~programowy prefetching. Stała
\kod{PREFETCH\_DIST}~=~8.
\begin{itemize}
\item \textbf{Prefetch wyprzedzający wiersz} \kod{A[(i+8)*cols]} ukrywa
  opóźnienie sprowadzenia kolejnego wiersza z~L3/pamięci.
\item \textbf{Prefetch wewnątrz wiersza} pod adresem \kod{+64}~bajty nakłada
  sprowadzanie danych na bieżące obliczenia.
\item \textbf{Odczyty wyrównane} (\kod{\_mm256\_load\_pd}) zamiast
  niewyrównanych --- przy założeniu wyrównania bufora do 32~bajtów.
\end{itemize}

\subsection{Wariant blokowany na cache (\kod{gemv\_avx\_fma\_blocked})}
Trzy warianty mikrojądra (L1/L2/L3) wybierane na podstawie rozmiaru zbioru
roboczego. Progi: \kod{L1\_THRESHOLD}~=~24~KB (75\,\% z~32~KB), \kod{L2\_THRESHOLD}
=~384~KB (75\,\% z~512~KB). Bloki dobrano w~osobnym etapie strojenia
(\kod{tuning.c}): L1~i~L3 --- \mbox{$64\times512$}, L2 --- \mbox{$128\times256$}.
\begin{itemize}
\item \textbf{Dyspozytor wg~rozmiaru} dobiera blok i~politykę prefetchingu do
  poziomu cache mieszczącego zbiór roboczy; jądro L1 nie używa prefetchingu
  (dane już w~cache).
\item \textbf{Margines 75\,\%} pozostawia miejsce na wektor $x$, wektor wynikowy
  i~dane wypierane przez system.
\item \textbf{Strojenie offline} --- konkretne rozmiary bloków są wynikiem
  przeszukania siatki parametrów programem \kod{tuning.c}, nie zaś analitycznego
  oszacowania\todocite{gemv-paper}.
\item \textbf{Makro \kod{DEFINE\_BLOCKED\_KERNEL}} generuje trzy warianty jądra
  z~jednego szablonu, eliminując powielanie kodu.
\end{itemize}

\subsection{Warianty z~mikrojądrem 4-wierszowym (\kod{gemv\_avx\_fma\_v2}, \kod{\_v3})}
Mikrojądro przetwarza jednocześnie 4~wiersze macierzy, mnożąc je przez wspólnie
wczytany fragment wektora $x$. Wariant \kod{v2} używa jednego akumulatora na
wiersz; wariant \kod{v3} --- dwóch akumulatorów na wiersz i~rozwinięcia o~16~kolumn.
\begin{itemize}
\item \textbf{Współdzielenie wektora $x$} przez 4~wiersze czterokrotnie zwiększa
  intensywność arytmetyczną względem przetwarzania wiersz po wierszu (jeden
  odczyt $x_j$ obsługuje cztery operacje FMA).
\item \textbf{Podwójne akumulatory w~\kod{v3}} (8~akumulatorów: 2~na wiersz)
  rozdzielają zależności FMA przy rozwinięciu o~16~kolumn, kosztem większej
  liczby zajętych rejestrów.
\item \textbf{Jedna redukcja pozioma na wiersz} pozostaje mimo rozwinięcia, dzięki
  scaleniu akumulatorów częściowych dopiero przed redukcją.
\end{itemize}

\subsection{Wariant wielowątkowy (\kod{gemv\_avx\_fma\_v3\_omp})}
Wielowątkowa wersja \kod{v3} z~progami uruchamiania zrównoleglenia.
\begin{itemize}
\item \textbf{Progi.} \kod{MIN\_ROWS\_FOR\_PARALLEL}~=~128,
  \kod{MIN\_ELEMENTS\_FOR\_PARALLEL}~=~$256\times256$,
  \kod{MAX\_ELEMENTS\_FOR\_PARALLEL}~=~8~M. Poniżej dolnych progów dominuje
  narzut tworzenia wątków; powyżej górnego progu problem jest ograniczony
  przepustowością pamięci i~wątki rywalizują o~kontroler pamięci, więc
  zrównoleglenie nie przynosi korzyści\todocite{gemv-paper}.
\item \textbf{Pojedynczy region \kod{\#pragma omp parallel}} z~ręcznym podziałem
  wierszy zamiast \kod{omp parallel for} --- eliminuje narzut harmonogramowania
  na iterację.
\item \textbf{Granice wątków wyrównane do 4~wierszy}, aby każdy fragment używał
  4-wierszowego mikrojądra SIMD.
\end{itemize}

\section{Mnożenie macierz--macierz (\kod{gemm})}
\label{sec:gemm}

Operacja \kod{gemm} realizuje
\begin{equation}
C \leftarrow \alpha\,A\,B + \beta\,C,\qquad
A\in\mathbb{R}^{M\times K},\ B\in\mathbb{R}^{K\times N},\ C\in\mathbb{R}^{M\times N}.
\label{eq:gemm}
\end{equation}
W~odróżnieniu od poprzednich jąder \kod{gemm} jest ograniczona mocą obliczeniową
(intensywność arytmetyczna rośnie z~rozmiarem), więc kluczowe stają się: kształt
mikrojądra maksymalizujący wykorzystanie jednostek FMA, pakowanie danych
zapewniające sekwencyjny dostęp oraz wielopoziomowe blokowanie na cache. Struktura
pętli i~pakowania jest wzorowana na kanonicznej ,,anatomii'' wysokowydajnego
\kod{gemm}\todocite{goto2008}.

\subsection{Wariant naiwny (\kod{gemm\_naive})}
Trzy zagnieżdżone pętle w~kolejności \kod{i,j,k} z~akumulacją iloczynu
wewnętrznego. Złożoność $O(N^3)$ przy bardzo słabej lokalności (kolumnowy odczyt
macierzy $B$).

\subsection{Zmiana kolejności pętli (\kod{gemm\_ikj})}
Kolejność \kod{i,k,j} z~wyniesieniem $\alpha A_{ik}$ przed pętlę po~$j$. Wiersz
$B_{k,:}$ jest wówczas odczytywany sekwencyjnie.
\begin{itemize}
\item \textbf{Sekwencyjny dostęp do $B$} poprawia wykorzystanie linii cache
  względem wariantu naiwnego.
\item \textbf{Wyniesienie mnożenia przez $\alpha$} --- raz na iterację~$k$ zamiast
  raz na parę $(i,j)$.
\end{itemize}

\subsection{Blokowanie trójpoziomowe (\kod{gemm\_blocked})}
Blokowanie z~rozmiarami \kod{BM}~=~64, \kod{BN}~=~64, \kod{BK}~=~256 (kolejność
pętli \kod{jc,pc,ic} a~następnie \kod{i,k,j}).
\begin{itemize}
\item \textbf{Blok $A$ o~rozmiarze $64\times256$} ($\approx$~128~KiB) i~analogiczny
  blok $B$ projektowane są pod rezydencję w~L2.
\item \textbf{Brak pakowania} --- wariant pokazuje zysk z~samego blokowania,
  jeszcze bez kosztu i~korzyści przepakowania danych.
\end{itemize}

\subsection{Pakowane mikrojądro $6\times8$ (\kod{gemm\_packed}, \kod{gemm\_omp})}
Pierwsza implementacja w~pełni zgodna ze schematem GotoBLAS/BLIS: dane $A$~i~$B$
są przepakowywane do paneli o~sekwencyjnym układzie, a~obliczenia wykonuje
mikrojądro $6\times8$ (Listing~\ref{lst:ukr6x8}). Stałe:
\kod{MR}~=~6, \kod{NR}~=~8, \kod{MC}~=~72, \kod{KC}~=~256, \kod{NC}~=~4080.
\begin{itemize}
\item \textbf{Kafelek $6\times8$} wykorzystuje 12~akumulatorów (6~wierszy
  $\times$~2~pasy po~4~liczby \kod{double}) plus 1~rozgłoszenie $A$ i~2~wektory $B$
  --- razem 15~z~16~rejestrów \kod{YMM}; jest to kanoniczny kafelek BLIS dla
  architektur Haswell/Zen\todocite{blis2015}.
\item \textbf{Układ pakowania.} Panel $A$ ma rozstaw \kod{MR} między wierszami
  jednej kolumny~$k$, co pozwala mikrojądru pobierać $a_i$ instrukcją
  \kod{vbroadcastsd} z~pamięci ciągłej; panel $B$ ma rozstaw \kod{NR}\todocite{goto2008}.
\item \textbf{Epilog \kod{vfmadd213pd}} liczy $C \mathrel{+}= \alpha\cdot\mathrm{acc}$
  w~miejscu, łącząc mnożenie przez $\alpha$ z~dodaniem starej wartości $C$.
\item \textbf{Obsługa krawędzi} przez skalarne jądro \kod{ukr\_6x8\_edge} dla
  niepełnych kafelków.
\item \textbf{Wariant wielowątkowy (\kod{gemm\_omp})} dzieli wymiar $M$ między
  wątki (zaokrąglając do \kod{MR}); każdy wątek pakuje własną kopię $B$ ---
  jest to najprostsza, lecz pamięciowo kosztowna strategia, stanowiąca punkt
  odniesienia dla wariantów ze współdzielonym $B$.
\end{itemize}

\begin{lstfloat}[htbp]
\lstset{language=C}
\begin{lstlisting}[caption={Pętla $K$ mikrojądra $6\times8$ z~12~akumulatorami (\kod{gemm\_packed.c}).}, label={lst:ukr6x8}]
for (int k = 0; k < kc; k++) {
    __m256d b0 = _mm256_load_pd(&B_pack[k*NR + 0]);
    __m256d b1 = _mm256_load_pd(&B_pack[k*NR + 4]);
    __m256d va;
    va = _mm256_broadcast_sd(&A_pack[k*MR + 0]);
    c00 = _mm256_fmadd_pd(va, b0, c00); c01 = _mm256_fmadd_pd(va, b1, c01);
    /* ... wiersze 1..4 ... */
    va = _mm256_broadcast_sd(&A_pack[k*MR + 5]);
    c50 = _mm256_fmadd_pd(va, b0, c50); c51 = _mm256_fmadd_pd(va, b1, c51);
}
\end{lstlisting}
\end{lstfloat}

\subsection{Pakowane mikrojądro $4\times12$ w~asemblerze (\kod{gemm\_zen3})}
\label{ssec:zen3}
Mikrojądro dostrojone specjalnie do Zen~3, z~pętlą~$K$ zapisaną we wstawce
asemblerowej (Listing~\ref{lst:ukr4x12}). Stałe:
\kod{MR\_Z}~=~4, \kod{NR\_Z}~=~12, \kod{MC\_Z}~=~192, \kod{KC\_Z}~=~240,
\kod{NC\_Z}~=~4080.
\begin{itemize}
\item \textbf{Kafelek $4\times12$.} 12~akumulatorów (4~wiersze $\times$~3~pasy)
  $+$~1~rozgłoszenie $A$ $+$~3~wektory $B$ wykorzystuje \emph{dokładnie}
  16~rejestrów \kod{YMM}. Analiza zliczania $\mu$-operacji wskazuje, że ten
  kształt daje korzystniejszy stosunek odczytów do operacji FMA niż kanoniczny
  $6\times8$ przy pełnym wysyceniu pliku rejestrów\todocite{agner-instr}.
\item \textbf{Wstawka asemblerowa} jest konieczna, ponieważ przy pełnym wysyceniu
  16~rejestrów kompilator GCC wprowadza zmienne tymczasowe i~przepełnia
  (\emph{spilluje}) dwa akumulatory na stos; jawne przypięcie akumulatorów do
  \kod{ymm4}--\kod{ymm15} eliminuje ten problem (szczegóły w~dokumencie
  \kod{INLINE\_ASM\_RATIONALE.md})\todocite{agner-microarch}.
\item \textbf{Brak programowego prefetchingu w~pętli~$K$.} Plik rejestrów jest
  pełny, a~obliczanie adresu prefetcha wymusiłoby spill akumulatorów; sprzętowy
  prefetcher strumieniowy Zen~3 jest dla odczytów sekwencyjnych wystarczający i~zalecany\todocite{agner-microarch}.
\item \textbf{Dobór faktorów blokowania.} \kod{KC\_Z}~=~240 dobrano tak, aby
  mikropanel $A$ ($\approx$~7,5~KiB) mieścił się w~L1d wraz z~rezydentnym
  fragmentem $B$; \kod{MC\_Z}~=~192 wynika z~analitycznego ograniczenia
  rezydencji panelu w~L2\todocite{goto2008}; \kod{NC\_Z}~=~4080 (panel
  $\approx$~7,5~MiB) odpowiada rezydencji w~L3 współdzielonej między rdzeniami
  CCX\todocite{amd-sog}.
\item \textbf{Pakery wektorowe} (plik \kod{gemm\_zen3\_pack.h}) realizują
  transpozycję bloków $4\times4$ macierzy $A$ instrukcjami
  \kod{unpack}/\kod{permute2f128} oraz trzy ciągłe kopie 256-bitowe na wiersz
  $B$; nagłówek jest współdzielony przez kilka jednostek translacji, aby każda
  otrzymała własną wersję inline.
\end{itemize}

\begin{lstfloat}[htbp]
\begin{lstlisting}[caption={Fragment pętli $K$ mikrojądra $4\times12$ we wstawce asemblerowej; akumulatory są przypięte do \kod{ymm4}--\kod{ymm15} (\kod{gemm\_zen3.c}).}, label={lst:ukr4x12}]
".balign 32                       \n"
"1:                               \n"
"vmovapd   0(%[B]), %%ymm1        \n"   // 3 wektory B (12 kolumn)
"vmovapd  32(%[B]), %%ymm2        \n"
"vmovapd  64(%[B]), %%ymm3        \n"
"vbroadcastsd 0(%[A]), %%ymm0     \n"   // rozgloszenie A[0]
"vfmadd231pd %%ymm1, %%ymm0, %%ymm4   \n"
"vfmadd231pd %%ymm2, %%ymm0, %%ymm5   \n"
"vfmadd231pd %%ymm3, %%ymm0, %%ymm6   \n"
/* ... wiersze 1..3 -> ymm7..ymm15 ... */
"addq $32, %[A]                   \n"
"addq $96, %[B]                   \n"
"decq %[k]                        \n"
"jnz 1b                           \n"
\end{lstlisting}
\end{lstfloat}

\subsection{Ścieżka małych macierzy (\kod{gemm\_zen3\_tiny})}
Dla \mbox{$\min(M,N,K)\le 96$} mikrojądro $4\times12$ działa bezpośrednio na
wejściowych macierzach wierszowych, bez pakowania i~bez \kod{aligned\_alloc}.
\begin{itemize}
\item \textbf{Pominięcie trzech narzutów} (alokacja, pakowanie, osobny przebieg
  \kod{apply\_beta}), które dominują przy rozmiarach rzędu mikrosekund; wzorowane
  na szybkich ścieżkach małych macierzy w~BLIS\todocite{blis2015}.
\item \textbf{Dodatkowe jądro $4\times4$} wektoryzuje częsty przypadek
  $N\bmod 12\in\{4,8\}$ (np. $N=64$, $N=256$), zastępując kosztowną pętlę
  skalarną.
\item \textbf{Epilog ,,tylko zapis''} dla $\beta=0$ pomija odczyt $C$.
\end{itemize}

\subsection{Współdzielony panel $B$ (\kod{gemm\_zen3\_par})}
Edukacyjny wariant wielowątkowy izolujący jeden zysk: współdzielenie panelu $B$.
\begin{itemize}
\item \textbf{Jeden wspólny \kod{B\_pack}} dla wszystkich wątków
  ($\approx$~7,5~MiB) zamiast kopii na wątek, dzięki czemu panel mieści się w~L3
  raz, a~nie wielokrotnie; każdy wątek pakuje własny panel $A$.
\item \textbf{Zrównoleglenie pętli \kod{ic}} (styl wielopętlowy BLIS) ---
  pętle \kod{jc}~i~\kod{pc} pozostają sekwencyjne\todocite{blis2015}.
\end{itemize}

\subsection{Warianty pośrednie (\kod{gemm\_zen3\_tuned}, \kod{gemm\_zen3\_sched})}
\begin{itemize}
\item \textbf{\kod{gemm\_zen3\_tuned}} dodaje trwały (\emph{persistent}) bufor
  roboczy reużywany między wywołaniami oraz wektorowe pakery, lecz bez dyspozytora
  rozmiarów --- izoluje wpływ pojedynczej optymalizacji.
\item \textbf{\kod{gemm\_zen3\_sched}} stosuje jawne harmonogramowanie: wektor $B$
  dla iteracji $k{+}1$ jest wczytywany na końcu iteracji $k$, aby ukryć
  pięciotaktowe opóźnienie odczytu z~L1 pod operacjami FMA poprzedniej iteracji
  (rozwinięcie pętli~$K$ o~4).
\end{itemize}

\subsection{Dyspozytor produkcyjny (\kod{gemm\_zen3\_best})}
Łączy wszystkie skuteczne optymalizacje pod jednym dyspozytorem.
\begin{itemize}
\item \textbf{Trójdrożny wybór ścieżki.} (i)~$\min(M,N,K)\le 96$ --- ścieżka
  ,,tiny''; (ii)~dostatecznie wiele bloków \kod{ic} ($M \ge n_t\cdot\kod{MC\_Z}$)
  --- współdzielony $B$ z~równoległą pętlą \kod{ic}; (iii)~w~przeciwnym razie
  --- podział wymiaru $M$ między wątki (osobne $B$), gdyż przy małym $M$
  współdzielony $B$ pozostawiłby część wątków bezczynnych.
\item \textbf{Równoległe pakowanie $B$} włączane tylko, gdy liczba pasów jest
  dostatecznie duża, aby zamortyzować barierę OpenMP ($\approx$~30~$\mu$s przy
  6~wątkach).
\item \textbf{Trwały bufor roboczy} eliminuje koszt \kod{aligned\_alloc} na każde
  wywołanie\todocite{blis2015}.
\end{itemize}

\subsection{Algorytm Strassena (\kod{gemm\_strassen})}
Wariant wykorzystujący schemat Strassena--Winograda o~siedmiu mnożeniach na
rozkładzie ćwiartkowym $A=\bigl[\begin{smallmatrix}A_{11}&A_{12}\\A_{21}&A_{22}\end{smallmatrix}\bigr]$
(analogicznie $B$, $C$). Zamiast ośmiu mnożeń klasycznych wykonuje się siedem,
kosztem dodatkowych przejść dodawania/odejmowania macierzy.
\begin{itemize}
\item \textbf{Próg \kod{STRASSEN\_MIN}~=~2048.} Poniżej tego rozmiaru
  12,5\,\%~oszczędności mnożeń nie równoważy kosztu osiemnastu przejść
  \kod{mat\_add}/\kod{mat\_sub}/\kod{mat\_copy}; próg wyznaczono empirycznie na
  Ryzen~5~5600\todocite{schwartz2023}.
\item \textbf{Przypadek bazowy} (poniżej progu) deleguje do
  \kod{gemm\_zen3\_best\_omp}.
\item \textbf{Zrównoleglone przejścia macierzowe.} Wszystkie 18~przejść jest
  zrównoleglonych po wymiarze wierszy (\kod{\#pragma omp parallel for}); przy
  $4096^3$ koszt pojedynczego przejścia spada z~$\sim$6~ms (1~wątek) do
  $\sim$1~ms (6~rdzeni).
\item \textbf{Kopiowanie ćwiartek przez bufory $S_1/S_2$.} Mikrojądro
  \kod{gemm\_zen3} zakłada, że rozstaw wiersza równa się wymiarowi macierzy
  (brak parametrów \kod{lda}/\kod{ldb}/\kod{ldc}); surowe ćwiartki dziedziczą
  rozstaw macierzy rodzica, więc przed każdym wywołaniem są kopiowane do ciągłych
  buforów.
\end{itemize}

\section{Splot dwuwymiarowy (\kod{conv})}
\label{sec:conv}

Splot realizowany jest dla pojedynczej próbki, kroku~1, bez dopełnienia, w~typie
\kod{float32}, dla danych w~układzie NCHW:
\begin{equation}
Y_{o,i,j} = \sum_{c=0}^{C_{in}-1}\sum_{p=0}^{K_H-1}\sum_{q=0}^{K_W-1}
            X_{c,\,i+p,\,j+q}\; W_{o,c,p,q},
\label{eq:conv}
\end{equation}
gdzie $X\in\mathbb{R}^{C_{in}\times H\times W}$,
$W\in\mathbb{R}^{C_{out}\times C_{in}\times K_H\times K_W}$,
$O_H=H-K_H+1$, $O_W=W-K_W+1$. Charakter obliczeń silnie zależy od kształtu jądra,
dlatego --- inaczej niż w~\kod{gemm} --- istnieje kilka jakościowo różnych ścieżek
(splot bezpośredni, $1\times1$ jako \kod{gemm}, Winograd) wybieranych przez
dyspozytor.

\subsection{Warianty naiwny, z~kolejnością pętli i~blokowany}
\begin{itemize}
\item \textbf{\kod{conv\_naive}} --- sześć zagnieżdżonych pętli, implementacja
  referencyjna o~słabej lokalności.
\item \textbf{\kod{conv\_reorder}} --- kolejność \kod{oc,ic,kh,kw,oh,ow};
  dla ustalonego $(oc,ic,kh,kw)$ wewnętrzne pętle czytają $X$ i~zapisują $Y$
  sekwencyjnie po~$ow$, a~waga $W_{oc,ic,kh,kw}$ jest niezmiennikiem wyniesionym
  przed pętlę.
\item \textbf{\kod{conv\_blocked}} --- blokowanie z~\kod{BLOCK\_OC}~=~32,
  \kod{BLOCK\_OH}~=~16, \kod{BLOCK\_IC}~=~32, ograniczające zbiór roboczy do L1/L2.
\end{itemize}

\subsection{Splot pakowany (\kod{conv\_packed}, \kod{conv\_omp})}
Mikrojądro o~kafelku \kod{CONV\_MR}~=~6 pozycji przestrzennych
$\times$~\kod{CONV\_NR}~=~16 kanałów wyjściowych. Wagi są przepakowywane do układu
z~$C_{out}$ jako wymiarem najbardziej wewnętrznym.
\begin{itemize}
\item \textbf{Kafelek $6\times16$.} 6~akumulatorów (po jednym na pozycję $ow$,
  każdy jako 2~pasy po~8~liczb \kod{float}) $+$~1~wektor wag $+$~1~rozgłoszona
  wartość $X$ --- razem 8~aktywnych rejestrów \kod{YMM}, z~zapasem na zmianę
  kolejności wykonania przez procesor.
\item \textbf{Przepakowanie wag} (\kod{conv\_pack\_W\_oc\_innermost}) tak, aby
  wartości dla kolejnych kanałów wyjściowych leżały ciągle --- mikrojądro
  rozgłasza wówczas wagę dla 16~kanałów jedną instrukcją.
\item \textbf{Wariant \kod{conv\_omp}} pakuje wagi raz (wspólny \kod{W\_pack}),
  a~następnie dzieli $C_{out}$ między wątki (zaokrąglając do \kod{CONV\_MR});
  każdy wątek pakuje własny fragment $X$.
\end{itemize}

\subsection{Splot $1\times1$ jako \kod{gemm} (\kod{conv\_1x1})}
Dla jąder $K_H=K_W=1$ splot jest czystym mnożeniem macierzy
$Y(C_{out},O_H O_W)=W(C_{out},C_{in})\cdot X(C_{in},O_H O_W)$.
\begin{itemize}
\item \textbf{Brak im2col.} W~układzie NCHW macierz $X$ jest już ułożona jako
  $(C_{in},O_H O_W)$, więc wywołanie kierowane jest bezpośrednio do
  \kod{cblas\_sgemm}, z~pominięciem bufora pośredniego.
\item \textbf{Fallback} do ścieżki pakowanej dla jąder innych niż $1\times1$.
\end{itemize}

\subsection{Splot Winograda $F(2\times2,3\times3)$ (\kod{conv\_winograd})}
Dla jąder $3\times3$ o~kroku~1 zastosowano algorytm Winograda, który redukuje
liczbę mnożeń na kafelek wyjściowy $2\times2$ z~36 (splot bezpośredni) do
16\todocite{lavin2016}. Wynik kafelka dany jest wzorem
\begin{equation}
Y = A^{\mathsf T}\Bigl[\,(G\,g\,G^{\mathsf T})\odot(B^{\mathsf T} d\,B)\,\Bigr] A,
\label{eq:winograd}
\end{equation}
gdzie $\odot$ oznacza mnożenie po współrzędnych, $g$ to filtr $3\times3$, $d$ ---
kafelek wejściowy $4\times4$, a~macierze transformacji to
\[
G=\begin{bmatrix}1&0&0\\[2pt]\tfrac12&\tfrac12&\tfrac12\\[2pt]\tfrac12&-\tfrac12&\tfrac12\\[2pt]0&0&1\end{bmatrix},\quad
B^{\mathsf T}=\begin{bmatrix}1&0&-1&0\\0&1&1&0\\0&-1&1&0\\0&1&0&-1\end{bmatrix},\quad
A^{\mathsf T}=\begin{bmatrix}1&1&1&0\\0&1&-1&-1\end{bmatrix}.
\]
Obliczenia rozbito na cztery etapy (Listing~\ref{lst:winograd}): transformacja
filtrów $U$, transformacja wejścia $V$, 16~mnożeń macierzy
$M_i=U_i\cdot V_i$ oraz transformacja wyjścia.
\begin{itemize}
\item \textbf{Redukcja mnożeń $36\to16$} ($2,25\times$) jest źródłem zysku;
  koszt to cztery transformacje kafelków $4\times4$ oraz 16~wywołań
  \kod{sgemm}\todocite{lavin2016}.
\item \textbf{Warunki stosowalności.} $K_H=K_W=3$, $O_H O_W \ge 16$ (amortyzacja
  narzutu transformacji) oraz parzyste $O_H$ i~$O_W$ (kafelkowanie $2\times2$).
\item \textbf{Macierze o~wpisach całkowitych/wymiernych} pozwalają zastąpić mnożenia
  w~transformacjach $B^{\mathsf T}dB$ i~$A^{\mathsf T}MA$ samymi dodawaniami
  i~odejmowaniami\todocite{castello2023}.
\item \textbf{Obsługa krawędzi.} Przy nieparzystych $O_H$/$O_W$ ostatni wiersz
  i~kolumna wyniku liczone są bezpośrednio (ułamek pracy rzędu $1/O_H$).
\end{itemize}

\begin{lstfloat}[htbp]
\lstset{language=C}
\begin{lstlisting}[caption={Transformacja wejścia Winograda $v=B^{\mathsf T} d\,B$ rozpisana na dodawania (\kod{conv\_winograd.c}).}, label={lst:winograd}]
for (int j = 0; j < 4; j++) {        // t = B^T d
    t[0][j] =  d[0][j] - d[2][j];
    t[1][j] =  d[1][j] + d[2][j];
    t[2][j] = -d[1][j] + d[2][j];
    t[3][j] =  d[1][j] - d[3][j];
}
for (int i = 0; i < 4; i++) {        // v = t B
    v[i][0] =  t[i][0] - t[i][2];
    v[i][1] =  t[i][1] + t[i][2];
    v[i][2] = -t[i][1] + t[i][2];
    v[i][3] =  t[i][1] - t[i][3];
}
\end{lstlisting}
\end{lstfloat}

\subsection{Splot blokowany kanałowo NCHWc (\kod{conv\_nchwc})}
Bezpośredni splot na układzie blokowanym kanałowo z~\kod{CB}~=~8 (jeden pas
\kod{YMM} w~\kod{float32}): $X$ przekształcane jest do
$(C_{in}/C_b,H,W,C_b)$, a~wagi do $(C_{out}/C_b,C_{in}/C_b,K_H,K_W,C_b,C_b)$.
\begin{itemize}
\item \textbf{\kod{CB}~=~8} odpowiada szerokości jednego pasa \kod{YMM} dla
  \kod{float}; mikrojądro wczytuje jeden wektor wag (8~kanałów wyjściowych) i~rozgłasza
  6~wartości przestrzennych $X$, akumulując do 6~rejestrów --- 8~aktywnych
  \kod{YMM}\todocite{onednn}.
\item \textbf{Brak rozgłaszania po stronie wag} --- rozgłaszane są wartości $X$,
  co eliminuje dodatkowe instrukcje na pętli wag.
\item \textbf{Fallback} do ścieżki pakowanej, gdy $C_{in}$ lub $C_{out}$ nie jest
  wielokrotnością~8.
\end{itemize}

\subsection{Dyspozytor splotu (\kod{conv\_zen3})}
Wybiera ścieżkę specjalizowaną na podstawie kształtu jądra, analogicznie do
\kod{gemm\_zen3\_best}:
\begin{itemize}
\item $K_H=K_W=1 \Rightarrow$ \kod{conv\_1x1} (SGEMM);
\item $K_H=K_W=3$ oraz $O_H O_W\ge16$ i~parzyste $O_H,O_W \Rightarrow$
  \kod{conv\_winograd};
\item w~przeciwnym razie $\Rightarrow$ \kod{conv\_nchwc}.
\end{itemize}
Region OpenMP znajduje się wewnątrz funkcji docelowych (w~transformacjach układu
i~pętlach głównych), a~nie wokół samego dyspozytora.

\subsection{Ścieżki im2col oraz biblioteki wzorcowe}
\begin{itemize}
\item \textbf{\kod{conv\_im2col\_openblas} / \kod{conv\_im2col\_blis}} ---
  rozwijają $X$ do macierzy kolumnowej (\kod{im2col}) i~wywołują SGEMM. Układ
  wierszowy $Y=W\cdot\mathrm{col}$ uzyskuje się sztuczką transpozycji: w~ujęciu
  kolumnowym (Fortran) operandy są oglądane jako $(N,K)$ i~$(K,C_{out})$, co daje
  ten sam układ bajtów co pożądane mnożenie wierszowe.
\item \textbf{\kod{conv\_libxsmm}} --- cienka nakładka na JIT-owane SGEMM
  biblioteki libxsmm, ładowanej dynamicznie; przy braku biblioteki funkcja jest
  pusta (pozycja w~benchmarku oznaczana jako N/A). Wykorzystano publiczne wejście
  \kod{libxsmm\_sgemm} ze względu na stabilność API między wersjami\todocite{libxsmm2016}.
\end{itemize}

\section{Zestawienie stałych projektowych}
\label{sec:stale}

Tabela~\ref{tab:stale} zestawia kluczowe stałe kafelków i~blokowania
poszczególnych jąder wraz z~ich uzasadnieniem sprzętowym.

\begin{table}[t!]
\begin{center}
\caption{Zestawienie kluczowych stałych projektowych jąder}\label{tab:stale}
\begin{tabular}{l|l|l}
Jądro / wariant & Stałe & Uzasadnienie \\
\hline
\kod{dot\_simd\_multiacc} & 8~akum., unroll~32 & ukrycie opóźnienia FMA \\
\kod{gemv} blokowany & progi 24~KB / 384~KB & 75\,\% pojemności L1 / L2 \\
\kod{gemm} $6\times8$ & MR=6, NR=8, MC=72, & kanoniczny kafelek BLIS, \\
                      & KC=256, NC=4080      & rezydencja L1/L2/L3 \\
\kod{gemm} $4\times12$ & MR=4, NR=12, MC=192, & pełne wysycenie 16~YMM, \\
                       & KC=240, NC=4080      & dostrojenie do Zen~3 \\
\kod{gemm\_strassen} & STRASSEN\_MIN=2048 & empiryczny próg opłacalności \\
\kod{conv} pakowany & CONV\_MR=6, CONV\_NR=16 & 8~aktywnych YMM \\
\kod{conv\_nchwc} & CB=8 & jeden pas YMM (float32) \\
\kod{conv\_winograd} & $F(2\times2,3\times3)$ & redukcja mnożeń $36\to16$ \\
\end{tabular}
\end{center}
\end{table}

% =====================================================================
% UWAGA: poniżej NIE umieszczono bibliografii. Znaczniki [CYT:klucz]
% w tekście wskazują miejsca na przyszłe odnośniki. Mapowanie kluczy:
%   agner-microarch -> Agner Fog, "The microarchitecture of Intel, AMD
%                      and VIA CPUs" (rejestry YMM, prefetch, pipeline Zen 3)
%   agner-instr     -> Agner Fog, "Instruction tables" (latencja/przepustowość FMA)
%   amd-sog         -> AMD, "Software Optimization Guide for AMD Family 19h
%                      Processors" (hierarchia cache L1/L2/L3)
%   goto2008        -> Goto & van de Geijn, "Anatomy of high-performance matrix
%                      multiplication", ACM TOMS 2008 (pięciopętlowe blokowanie, pakowanie)
%   blis2015        -> Van Zee & van de Geijn, "BLIS: A framework...", ACM TOMS 2015
%                      (mikrojądro 6x8, shared-B, trwały workspace, ścieżka tiny)
%   schwartz2023    -> Schwartz & Vaknin, "Pebbling Game and Alternative Basis...",
%                      SIAM J. Sci. Comput. 2023 (próg crossover Strassena)
%   lavin2016       -> Lavin & Gray, "Fast Algorithms for Convolutional Neural
%                      Networks", CVPR 2016 (Winograd F(2,3), macierze transformacji)
%   castello2023    -> Castello i in., "Efficient and portable Winograd
%                      convolutions...", J. Supercomput. 2023 (przenośność AVX2)
%   libxsmm2016     -> Heinecke i in., "LIBXSMM: Accelerating Small Matrix
%                      Multiplications...", SC 2016
%   onednn          -> dokumentacja oneDNN / ZenDNN (układ NCHWc)
%   gemv-paper      -> własny artykuł autora o GEMV na Zen 3 (przepustowości, strojenie)
% =====================================================================

\end{document}
